% !TeX root = ../navier-stokes-manuscript.tex
\clearpage
\section{The endpoint burden is one rectangle}\label{sec:TheMissingEstimateAndGlobalSmoothness}

The full rectangle family is not needed to rule out a finite last time.
One bound would do:
\[
  \sup_{0<T<T_*}\mathfrak R_{0,1}(u;T)<\infty.
\]
Its two adjacent rows put \(u\) in \(C([0,T_*];H^1)\), hence in
\(L^\infty(I_*;L^3)\).  A finite maximal time requires that same \(L^3\) norm to
escape.  The two conclusions cannot coexist.

That rectangle bound has not been proved from the datum.  The result below is therefore
a reduction, not a closure: if \(T_*\) is finite, then \(\mathfrak R_{0,1}\) must diverge
as the cutoff reaches \(T_*\).  The higher rectangles answer a separate question.  If
the missing family were available, they would carry the entire mixed velocity--pressure
history, not merely the first continuation norm, to the endpoint.

\subsection{The lowest rectangle meets the endpoint criterion}\label{sub:FromFiniteRectanglesToClassicalEstimates}

Suppose \(T_*<\infty\), and suppose
\begin{equation}\label{eq:LowRectangleConditionalBound}
  \sup_{0<T<T_*}\mathfrak R_{0,1}(u;T)<\infty.
\end{equation}
Monotone exhaustion in the definition of \(\mathfrak R_{0,1}\) gives
\[
  u\in L^2(I_*;H^2),
  \qquad
  u_t\in L^2(I_*;L^2).
\]
The adjacent trace identity carries these two rows into the missing middle space:
\[
  u\in C([0,T_*];H^1).
\]
Since \(H^1(\R^3)\hookrightarrow L^3(\R^3)\),
\begin{equation}\label{eq:EndpointCriterionFromLowRectangle}
  \sup_{0\leq t\leq T_*}\norm{u(t)}_3<\infty.
\end{equation}
This is where the lowest rectangle meets the finite maximal-time criterion.
The endpoint theorem of Escauriaza, Seregin, and \v{S}ver\'ak says that a uniform
\(L^3\) bound continues the solution through \(T_*\)
\parencite{EscauriazaSereginSverak2003}.  Under the standing contradiction hypothesis,
\begin{equation}\label{eq:FiniteTimeLThreeEscape}
  T_*<\infty
  \quad\Longrightarrow\quad
  \sup_{0<t<T_*}\norm{u(t)}_3=\infty.
\end{equation}
Equation \eqref{eq:EndpointCriterionFromLowRectangle} gives the opposite conclusion.
Thus
\begin{equation}\label{eq:LowRectangleEndpointReduction}
  T_*<\infty
  \quad\Longrightarrow\quad
  \sup_{0<T<T_*}\mathfrak R_{0,1}(u;T)=\infty.
\end{equation}
This is the exact remaining burden.  A datum-side proof of
\eqref{eq:LowRectangleConditionalBound} would rule out a finite maximal time, but no
such proof has been supplied.

\subsection{What the full rectangle family adds}\label{sub:FullRectangleFamilyAdds}

Now assume the stronger whole-terminal rectangle condition stated in
\Cref{thm:BodyWholeTerminalRectangleTheorem}.
Its \((0,1)\) member gives the contradiction above.
The remaining coordinates answer a different question: what does the entire mixed
history carry to the putative endpoint under that condition?

Fix
\(r,m\in\Nzero\), choose \(N\geq r\), and take the rectangle of middle height
\(M=m\).  Its adjacent rows give
\[
  V_r\in L^2(I_*;H^{m+1}),
  \qquad
  V_{r+1}=\partial_tV_r\in L^2(I_*;H^{m-1}).
\]
\Cref{lem:AdjacentHilbertScaleTrace} gives
\begin{equation}\label{eq:BodyUniformMixedSobolevCoordinate}
  \sup_{0\leq t\leq T_*}
  \norm{\partial_t^ru(t)}_{H^m}
  <\infty.
\end{equation}
Nothing in this step couples \(r\) to \(m\).  We may choose any finite temporal
response and then ask for any finite spatial height inside that same row.

Given \(k\in\Nzero\) and \(2\leq p\leq\infty\), choose an integer \(m\) with
\begin{equation}\label{eq:SobolevProjectionThreshold}
  m>k+\frac32-\frac3p.
\end{equation}
The spatial Sobolev embedding \parencite[Theorem~7.28]{Cerda2010} gives
\begin{equation}\label{eq:MasterSpatialProjection}
  \norm{\nabla^k\partial_t^ru(t)}_{L^p}
  \leq C_{k,m,p}\norm{\partial_t^ru(t)}_{H^m}.
\end{equation}
Combining this with \eqref{eq:BodyUniformMixedSobolevCoordinate},
\begin{equation}\label{eq:MasterUniformMixedEstimate}
  \nabla^k\partial_t^ru
  \in L^\infty(I_*;L^p(\R^3)).
\end{equation}
Since \(T_*\) is finite under the standing contradiction hypothesis,
\begin{equation}\label{eq:MasterFiniteTimeMixedEstimate}
  \nabla^k\partial_t^ru
  \in L^q(I_*;L^p(\R^3))
  \qquad(1\leq q\leq\infty).
\end{equation}
The pressure is carried by these same finite coordinates.
For fixed \(r,m\in\Nzero\), the normalized pressure formula \eqref{eq:BodyTemporalPressureRecurrence} contains a finite sum of products of the rows \(V_0,\ldots,V_r\).
The adjacent traces place those rows continuously in every finite Sobolev space, so Sobolev multiplication and the Riesz-transform bounds give
\begin{equation}\label{eq:MasterPressureSobolevCoordinate}
  \partial_t^rp=Q_r
  \in C([0,T_*];H^m(\R^3)).
\end{equation}
For finite \(k\), \(2\leq p\leq\infty\), and \(1\leq q\leq\infty\), choose \(m>k+\tfrac32-\tfrac3p\).
The same spatial embedding \parencite[Theorem~7.28]{Cerda2010} and the finiteness of \(T_*\) then give
\begin{equation}\label{eq:MasterFiniteTimePressureEstimate}
  \nabla^k\partial_t^rp
  \in L^q(I_*;L^p(\R^3)).
\end{equation}
Each requested norm uses one finite rectangle, and the compatible intersection keeps
these finite coordinates together without replacing them by an all-orders sum.

\subsection{The whole history reaches the endpoint}\label{sub:WholeHistoryReachesEndpoint}

The full family gives more than a list of norms.  It gives every row a terminal value.
Fix \(m\).  A rectangle with middle height \(m+1\) contains both
\[
  u\in L^2(I_*;H^m),
  \qquad
  u_t\in L^2(I_*;H^m).
\]
Thus \(u\in H^1(I_*;H^m)\) and has a limit \(u_*^{(m)}\in H^m\) as
\(t\uparrow T_*\).  If \(k>m\), the limits in \(H^k\) and \(H^m\) agree after the
embedding \(H^k\hookrightarrow H^m\), because both are limits of the same \(u(t)\).
The compatible limits therefore define one endpoint
\begin{equation}\label{eq:BodyCommonSmoothEndpoint}
  u_*:=\lim_{t\uparrow T_*}u(t)
  \in\Hinfty_\sigma(\R^3).
\end{equation}

The same argument applies to every temporal row.  For finite \(N,m\), choose a
rectangle containing both \(V_N\) and \(V_{N+1}\) in \(L^2(I_*;H^m)\).  Then
\(V_N\in H^1(I_*;H^m)\) and has a terminal value \(V_N^*\).  Compatibility in
\(m\) places each \(V_N^*\) in \(\Hinfty\).  Passing to the endpoint in the
normalized recurrences gives
\begin{align}
  Q_N^*
  &=R_iR_j\sum_{a=0}^{N}\binom{N}{a}V_{a,i}^*V_{N-a,j}^*,
  \label{eq:BodyEndpointPressureRecurrence}
  \\
  V_{N+1}^*
  &=\nu\Delta V_N^*
  -\sum_{a=0}^{N}\binom{N}{a}(V_a^*\cdot\nabla)V_{N-a}^*
  -\nabla Q_N^*.
  \label{eq:BodyEndpointVelocityRecurrence}
\end{align}
The left-hand history has reached \(T_*\) with the velocity and pressure jet generated
by \(u_*\), not merely with a collection of time-integrated bounds.

Smooth local existence from \(u_*\) now produces a solution to the right of \(T_*\).
Its jet is generated by the same recurrences, and strong uniqueness joins it to the
history on the left.  This is a second continuation route.  It identifies the complete
endpoint jet, whereas the first route needs only the \(L^3\) bound.  It also spends the
whole rectangle family, and so remains conditional on that still-unproved family.

\subsection{The critical-height roof}\label{sub:ClassicalContinuationCriteriaInsideFiniteRectangles}

Under the same rectangle condition, the \(H^1\) trace also puts the original physical
picture back under a finite roof.  Since
\(H=\tfrac12\norm{\Lambda^{1/2}u}_2^2\),
\begin{equation}\label{eq:BodyCriticalHeightRoof}
  H_{\max}
  :=\sup_{0\leq t\leq T_*}H(t)
  \leq\frac12\sup_{0\leq t\leq T_*}\norm{u(t)}_{H^1}^2
  <\infty.
\end{equation}
Fix \(H_0>0\), set \(L_k:=2^kH_0\) for \(k\in\Nzero\), and let
\[
  N_{\mathrm{levels}}
  :=
  \#\bigl\{k\in\Nzero:\text{there is a }t\in[0,T_*]\text{ with }H(t)=L_k\bigr\}.
\]
Only finitely many of these dyadic levels can be reached beneath the roof, and
\begin{equation}\label{eq:BodyFiniteZenoOctaveCount}
  N_{\mathrm{levels}}
  \leq
  \begin{cases}
    0,&H_{\max}<H_0,\\[0.25em]
    1+\left\lfloor\log_2(H_{\max}/H_0)\right\rfloor,
      &H_{\max}\geq H_0.
  \end{cases}
\end{equation}
No estimate on the duration of an individual crossing is needed.
The critical-height ladder is finite, which rules out the unbounded height sequence forced by a finite maximal time.
This alone does not rule out shrinking spatial widths; the two vortex readings below test those separate geometric claims.

The next low rectangle, \(\cR_{0,2}\), controls how the height moves beneath that
roof:
\begin{equation}\label{eq:LowRectangleZeroTwo}
  u\in L^2(I_*;H^3),
  \qquad
  u_t\in L^2(I_*;H^1),
  \qquad
  u\in C([0,T_*];H^2).
\end{equation}
Indeed,
\[
  H'(t)
  =\operatorname{Re}
  \pair{\Lambda^{1/2}u(t)}{\Lambda^{1/2}u_t(t)}_{L^2},
\]
and Cauchy--Schwarz in time gives
\begin{equation}\label{eq:CriticalHeightTotalVariationBound}
  \int_0^{T_*}|H'(t)|\,dt
  \leq
  \norm{\Lambda^{1/2}u}_{L^2_tL^2_x}
  \norm{\Lambda^{1/2}u_t}_{L^2_tL^2_x}
  <\infty.
\end{equation}
Also \(E_{3/2}\in L^1(I_*)\), so the first completed row
\[
  \mathcal P_{1/2}=H'+2\nu E_{3/2}
\]
shows that the signed critical production is integrable on \(I_*\).
Integrating that same row identifies the signed prefix exactly:
\begin{equation}\label{eq:BodySignedCriticalPrefix}
  \int_0^t\mathcal P_{1/2}(s)\,ds
  =H(t)-H(0)+2\nu\int_0^tE_{3/2}(s)\,ds.
\end{equation}
The critical-height roof and the finite \(E_{3/2}\) integral give
\begin{equation}\label{eq:BodySignedCriticalPrefixBound}
  A_*(T_*)
  :=\sup_{0<t<T_*}\int_0^t\mathcal P_{1/2}(s)\,ds
  <\infty.
\end{equation}
The scalar signed-prefix bound is now a consequence of the compatible intersection.
It was not used to produce the mixed jet, its finite rectangles, or their whole-terminal bounds.
The height can rise and fall sharply, but it cannot accumulate an infinite net rise before \(T_*\).

The same rectangle also gives the physical control needed for the vortex.  In three dimensions,
\(H^3\hookrightarrow W^{1,\infty}\), and hence
\[
  \int_0^{T_*}\norm{\nabla u(t)}_\infty\,dt
  \leq
  T_*^{1/2}\norm{u}_{L^2(I_*;H^3)}<\infty.
\]
Since \(\omega=\nabla\times u\) and \(S=\tfrac12(\nabla u+\nabla u^{\mathsf T})\),
we obtain the full finite stretching action
\begin{equation}\label{eq:FiniteLipschitzVorticityStrainAction}
  \int_0^{T_*}
  \left(
    \norm{\nabla u(t)}_\infty
    +\norm{\omega(t)}_\infty
    +\norm{S(t)}_\infty
  \right)dt
  <\infty.
\end{equation}

\subsection{The two vortex readings}\label{sub:ClosingCriticalHeightZenoAndTheVortex}

The mixed intersection lets us test two distinct readings of a narrowing core.
They are not an exhaustive classification of singular geometry.
They isolate two different claims a vortex picture can make: either the visible core follows fixed material labels, or its membership changes as the concentration moves through the fluid.

Begin with fixed material labels.
Let \(X(a,t)\) solve
\[
  \partial_tX(a,t)=u(X(a,t),t),
  \qquad
  X(a,0)=a.
\]
For two trajectories in that core and \(t<T_*\),
\[
  \frac d{dt}|X(a,t)-X(b,t)|
  \geq
  -\norm{\nabla u(t)}_\infty|X(a,t)-X(b,t)|.
\]
The finite action in \eqref{eq:FiniteLipschitzVorticityStrainAction} gives
\begin{equation}\label{eq:BodyMaterialCoreLowerBound}
  |X(a,t)-X(b,t)|
  \geq
  |a-b|
  \exp\!\left(
    -\int_0^t\norm{\nabla u(s)}_\infty\,ds
  \right)
  >0
\end{equation}
The same \(L^2(I_*;H^3)\) row gives \(u\in L^1(I_*;L^\infty)\), so each trajectory
extends continuously to \(T_*\).  Passing to the limit above gives
\[
  |X(a,T_*)-X(b,T_*)|
  \geq |a-b|\exp\!\left(-\int_0^{T_*}\norm{\nabla u(s)}_\infty\,ds\right)>0.
\]
A material core begins with distinct particles, and those particles remain distinct at
the terminal time.  The drawn radius may shrink, but it cannot represent the collapse of
that same labelled core to one point.

Keep this first reading and now put the vortex exactly on its centreline.
The finite coordinate \((r,m)=(0,3)\) in \eqref{eq:BodyUniformMixedSobolevCoordinate} gives
\begin{equation}\label{eq:BodyUniformLipschitzBound}
  L_*:=\sup_{0\leq t\leq T_*}\norm{\nabla u(t)}_\infty<\infty.
\end{equation}
At a chosen centre \(x_c(t)\),
\[
  |u(t,x)-u(t,x_c)|\leq L_*|x-x_c|.
\]
Let \(C_r(t)\) be a circle of radius \(r\) in a cross-section normal to the vortex axis.
The constant vector \(u(t,x_c)\) integrates to zero around the closed curve, so
\begin{equation}\label{eq:BodyCenteredCirculationBound}
  \left|
    \oint_{C_r(t)}u(t,x)\cdot d\ell
  \right|
  \leq2\pi L_*r^2
  \longrightarrow0
  \qquad(r\downarrow0).
\end{equation}
Perfect centring does not evade the Lipschitz bound.  If the radius tends to zero while
the velocity remains Lipschitz, the relative tangential speed is \(O(r)\) and the
circulation is \(O(r^2)\); no nonzero circulation survives at the limiting point.

The second reading lets the visible core recruit new material as earlier labels leave
it.
Particle separation no longer follows the changing membership of that core, so the test moves to the whole-field high-frequency tail.
For \(s>1/2\) and \(K\geq1\), define
\[
  H_{>K}(t)
  :=\frac12\int_{|\xi|>K}
  |\xi|\,|\widehat u(t,\xi)|^2\,d\xi.
\]
Then
\begin{equation}\label{eq:BodyUltravioletCriticalTail}
  H_{>K}(t)
  \leq
  \frac12K^{1-2s}\norm{u(t)}_{\dot H^s}^2.
\end{equation}
Choose an integer \(m\geq s\).
Equation \eqref{eq:BodyUniformMixedSobolevCoordinate} with \(r=0\), together with
\(H^m\hookrightarrow\dot H^s\), gives
\begin{equation}\label{eq:BodyUniformUltravioletVanishing}
  \lim_{K\to\infty}
  \sup_{0\leq t\leq T_*}H_{>K}(t)=0.
\end{equation}
Even when material is continually recruited, the complete velocity field cannot retain a fixed amount of critical content beyond every frequency.

These two readings now have separate conclusions.
Fixed material labels retain positive separation at \(T_*\), and a centred radius carries vanishing circulation as it shrinks.
Changing membership cannot preserve a nonzero whole-field critical tail beyond every frequency.

\subsection{Where the argument stops}\label{sub:NoFiniteSingularTime}

The proved conclusion is \eqref{eq:LowRectangleEndpointReduction}: if the maximal time
is finite, then \(\mathfrak R_{0,1}\) must become unbounded as the cutoff approaches
\(T_*\).  To conclude instead that \(T_*=\infty\), one must bound that same rectangle
from \(u_0\) and \(\nu\), with a constant that does not deteriorate as
\(T\uparrow T_*\).  No estimate in the preceding sections supplies that bound.

The stronger consequences above do not pay this debt.  They say what follows once the
whole rectangle family is assumed: the \(L^3\) contradiction, the common smooth
endpoint, the complete velocity--pressure jet, the same-history restart, the finite
critical-height ladder, and the two vortex conclusions.  None of those consequences
produces the missing datum-to-rectangle estimate.  The global argument remains open at
exactly that edge.

\section{What the finite rectangles see, and where they stop}\label{sec:FurtherConsequencesAndEulerBoundary}

The changing-membership vortex leaves one escape route to test.  Its activity can move
from shell to shell as the visible core recruits new fluid, even if no fixed
high-frequency tail stays large.  On every strict interval \(0<T<T_*\), the next finite
rectangle rules out an unlimited version of that migration: the time-integrated
vorticity strengths are summable over all dyadic shells.

This statement does not use the missing terminal estimate.  Ordinary smooth persistence
already makes each fixed rectangle finite when \(T<T_*\).  What remains open is whether
the relevant bound survives as \(T\uparrow T_*\).  Keeping that quantifier visible also
shows where the construction stops: viscosity supplies the rung that exposes the next
spatial height, but no argument here has yet bounded that rung uniformly from the datum.

\subsection{A finite action across every scale}\label{sub:AFiniteSumAcrossScales}

Fix \(T<T_*\).  The strict rectangle \(\cR_{0,2}[u_0;T]\) contains
\[
  u\in L^2(0,T;H^3).
\]
Let \((\Delta_j)_{j\geq-1}\) be an inhomogeneous Littlewood--Paley decomposition
\parencite[Chapter~2]{BahouriCheminDanchin2011}.  For \(j\geq0\), the block
\(\Delta_j\omega\) lies near wavelength \(2^{-j}\), and
\(\norm{\Delta_j\omega}_\infty\) records the strongest rotation in that shell.  Write
\[
  \norm{\omega}_{B^0_{\infty,1}}
  :=\sum_{j\geq-1}\norm{\Delta_j\omega}_\infty.
\]
Bernstein's inequality gives
\[
  \norm{\Delta_j\omega}_\infty
  \leq C2^{3j/2}\norm{\Delta_j\omega}_2
  \qquad(j\geq0).
\]
Insert \(2^{3j/2}=2^{-j/2}2^{2j}\), sum, and apply Cauchy--Schwarz in \(j\):
\begin{align*}
  \sum_{j\geq0}\norm{\Delta_j\omega}_\infty
  &\leq
  C\left(\sum_{j\geq0}2^{-j}\right)^{1/2}
  \left(\sum_{j\geq0}2^{4j}\norm{\Delta_j\omega}_2^2\right)^{1/2}
  \\
  &\leq C\norm{\omega}_{H^2}
  \leq C\norm{u}_{H^3}.
\end{align*}
The shell estimate costs \(3/2\) derivatives.  The \(H^2\) vorticity weight supplies
two, leaving the square-summable factor \(2^{-j/2}\).  After the low block is included,
\begin{equation}\label{eq:VorticityBesovFromHThree}
  \norm{\omega}_{B^0_{\infty,1}}
  \leq C\norm{u}_{H^3}.
\end{equation}
Cauchy--Schwarz in time now gives
\begin{equation}\label{eq:FiniteDyadicVorticityAction}
  \sum_{j\geq-1}
  \int_0^{T}\norm{\Delta_j\omega(t)}_\infty\,dt
  =\int_0^{T}\norm{\omega(t)}_{B^0_{\infty,1}}\,dt
  <\infty.
\end{equation}
The conclusion concerns the sum, not the number of active shells.  Activity may pass
through infinitely many scales, but its time-integrated strength cannot form a
nonsummable cascade before this fixed \(T\).  Nothing here supplies a bound uniform in
\(T\uparrow T_*\), an analytic radius, or a weighted sum over every derivative order.

\subsection{The same endpoint quantifier in classical and local form}\label{sub:ClassicalCriteriaInsideFiniteRectangles}

The shell sum is one reading of a strict finite rectangle.  The familiar continuation
criteria are three others.  On the same fixed interval, \(\cR_{0,1}[u_0;T]\) gives
\[
  u\in C([0,T];H^1),
  \qquad
  u\in L^2(0,T;H^2),
\]
and therefore
\begin{equation}\label{eq:StrictIntervalClassicalCriteria}
  \sup_{0\leq t\leq T}\norm{u(t)}_3<\infty,
  \qquad
  u\in L^2(0,T;L^\infty),
  \qquad
  \nabla u\in L^2(0,T;L^3).
\end{equation}
These are, respectively, the critical norm in the
Escauriaza--Seregin--\v{S}ver\'ak endpoint theorem
\parencite{EscauriazaSereginSverak2003}, the pair \((q,p)=(2,\infty)\) on the
Ladyzhenskaya--Prodi--Serrin line
\parencite{Prodi1959,Serrin1962,Serrin1963,Ladyzhenskaya1967}, and the gradient
criterion of Beir\~ao da Veiga \parencite{BeiraoDaVeiga1995}.  They are not three new
closure arguments.  They show that one low block reads velocity escape,
scale-balanced integrability, and vortex-producing shear at once.

The local singularity picture has the same boundary.  Smooth persistence of the
normalized velocity and pressure on \([0,T]\) gives
\begin{equation}\label{eq:FiniteIntervalVelocityPressureSupremum}
  K_T
  :=\sup_{0\leq t\leq T}
  \left(
    \norm{u(t)}_{L^\infty}
    +\norm{p(t)}_{L^\infty}
  \right)
  <\infty.
\end{equation}
For
\[
  Q_\rho(z_0)
  =B_\rho(x_0)\times(t_0-\rho^2,t_0)
  \subset\R^3\times(0,T],
\]
the scale-invariant velocity--pressure content satisfies
\begin{equation}\label{eq:VanishingLocalScaleInvariantContent}
  \rho^{-2}
  \iint_{Q_\rho(z_0)}
  \left(|u|^3+|p|^{3/2}\right)\,dx\,dt
  \leq
  C\rho^3\left(K_T^3+K_T^{3/2}\right)
  \longrightarrow0
  \qquad(\rho\downarrow0).
\end{equation}
Thus every point strictly before \(T_*\) loses the positive scale-invariant content
required by the Caffarelli--Kohn--Nirenberg singularity picture.  This does not remove a
candidate terminal face: the constants in \eqref{eq:StrictIntervalClassicalCriteria}
and \eqref{eq:FiniteIntervalVelocityPressureSupremum} may still deteriorate as
\(T\uparrow T_*\).  Section~4 isolates exactly the estimate that would prevent that
deterioration.

\Needspace{8\baselineskip}
\subsection{The ascent is viscous}\label{sub:WhyTheEulerTowerIsDifferent}

The fixed-interval bounds above might make it look as though writing the mixed temporal
jet is enough.  Euler has that jet too, so compare the two histories directly.  Write
\[
  V_n^{\mathrm{NS}}:=\partial_t^nu^{\mathrm{NS}},
  \qquad
  Q_n^{\mathrm{NS}}:=\partial_t^np^{\mathrm{NS}},
\]
and define \(V_n^{\mathrm E},Q_n^{\mathrm E}\) from a separate Euler solution.  Their
temporal recurrences are
\begin{align}
  V_{n+1}^{\mathrm{NS}}
  &+
  \sum_{a=0}^{n}\binom{n}{a}
  (V_a^{\mathrm{NS}}\cdot\nabla)V_{n-a}^{\mathrm{NS}}
  +\nabla Q_n^{\mathrm{NS}}
  =\nu\Delta V_n^{\mathrm{NS}},
  \label{eq:NSTemporalRecurrenceComparison}
  \\
  V_{n+1}^{\mathrm E}
  &+
  \sum_{a=0}^{n}\binom{n}{a}
  (V_a^{\mathrm E}\cdot\nabla)V_{n-a}^{\mathrm E}
  +\nabla Q_n^{\mathrm E}
  =0.
  \label{eq:EulerTemporalRecurrenceComparison}
\end{align}
Both retain differentiated transport and the nonlocal pressure response.  Only
Navier--Stokes retains \(\nu\Delta V_n^{\mathrm{NS}}\), the linear term that exposes two
additional spatial derivatives of row \(n\).

The energy rows show the same difference.  Set
\[
  E_s^{\mathrm{NS}}:=\frac12\norm{\Lambda^su^{\mathrm{NS}}}_2^2,
  \qquad
  E_s^{\mathrm E}:=\frac12\norm{\Lambda^su^{\mathrm E}}_2^2,
\]
with transfer terms \(\mathcal P_s^{\mathrm{NS}}\) and \(\mathcal P_s^{\mathrm E}\).
Then
\begin{equation}\label{eq:NSEulerEnergyComparison}
  (E_s^{\mathrm{NS}})'
  =\mathcal P_s^{\mathrm{NS}}-2\nu E_{s+1}^{\mathrm{NS}},
  \qquad
  (E_s^{\mathrm E})'
  =\mathcal P_s^{\mathrm E}.
\end{equation}
For Navier--Stokes, repeated differentiation and division by \((-2\nu)^n\) isolate
\(E_{n+1/2}^{\mathrm{NS}}\) in \eqref{eq:BodyCompletedCriticalRow}.  Euler keeps the
nonlinear mixing of time and space but has no signed viscous rung that isolates the next
Sobolev height.  This does not mean Euler lacks ordinary propagation of spatial
regularity.  It means that the completed-row ascent used here is no longer available.

An Euler history already known to lie in an all-orders compatible intersection would
satisfy every fixed Sobolev bound and, in particular, the Beale--Kato--Majda condition
\[
  \int_0^T\norm{\omega^{\mathrm E}(t)}_\infty\,dt<\infty
\]
\parencite{BealeKatoMajda1984}.  Nothing above constructs that Euler intersection, and
no constant in the fixed-viscosity argument is claimed to remain bounded as
\(\nu\downarrow0\).

The two boundaries are now separate.  For Navier--Stokes, every strict finite rectangle
is available, but the datum-side estimate uniform up to a putative \(T_*\) remains
unproved.  If the open \(\mathfrak R_{0,1}\) estimate is supplied, Section~4 rules out a
finite maximal time and the strict-interval conclusions above range over every finite
\(T\).  For Euler, the viscous rung used by this ascent is absent, and no inviscid
transfer has been obtained.  These calculations identify what the rectangle mechanism
sees and where it stops; they close neither gap.
